\documentclass[a4paper]{article}
\usepackage[margin=2.5cm]{geometry} % márgenes de article por defecto: ~4 cm
\usepackage{titling}
\setlength{\droptitle}{-2cm} % \maketitle deja mucho aire arriba del título
\usepackage[spanish]{babel}
\usepackage{amsmath}
\usepackage[md]{juliatex} % [md]: además analisis.md (pandoc)

\title{Análisis de Datos Científicos}
\author{Miguel Ferrera}

\begin{document}
\maketitle

\section{Introducción al modelo}

Analizamos una función trigonométrica amortiguada, definida por cada dato nuevo.
\begin{equation}
  f(x) = e^{-0.1x} \sin(x).
\end{equation}

\section{Simulación en Julia}

Definimos la función y evaluamos algunos puntos:

\begin{juliacode}
f(x) = exp(-0.1x) * sin(x)
x = 0:0.1:50
println("Puntos evaluados: ", length(x))
round.(f.(0:2:8), digits = 3)
\end{juliacode}

La gráfica se genera aparte, sin mostrar el código dependiente:

\begin{juliarun}
using Plots
plot(x, f.(x), label = "f(x)", lw = 2, xlabel = "Tiempo", ylabel = "Amplitud")
savefig("amortiguada.pdf")
savefig("amortiguada.svg") # para analisis.md (GitHub no muestra PDF)
\end{juliarun}

\begin{figure}[htbp]
  \centering
  \juliagraphics[width=0.75\linewidth]{amortiguada.pdf}
  \caption{Función amortiguada generada por Julia al compilar.}%
  \label{fig:amortiguada}
\end{figure}

Como se ve en la figura~\ref{fig:amortiguada}, la amplitud decrece
exponencialmente. El máximo calculado es \juliav{round(maximum(f.(x)), digits = 4)}.

\end{document}
